\chapter{Bản đồ thanh ghi CSR}
\label{app:csr}

Phụ lục này liệt kê đầy đủ bản đồ thanh ghi của khối \CsrRegisters{} (\CsrRegisterCount{} thanh ghi cùng vùng bảng
DAT \DatEntryCount{} mục), bổ trợ cho \cref{sec:rtl-csr}. Địa chỉ tính theo byte, độ rộng mỗi thanh ghi là \RegisterDataWidth{}-bit.

\begin{longtable}{L{4.6cm} C{2.0cm} L{6.4cm}}
  \toprule
  \textbf{Thanh ghi} & \textbf{Địa chỉ} & \textbf{Chức năng} \\
  \midrule
  \endhead
  \texttt{HC\_CONTROL}    & 0x004 & Điều khiển bộ điều khiển (\texttt{BUS\_ENABLE}, \texttt{IBA\_INCLUDE} bật phát \IThreeCBroadcastAddress{}, \texttt{ABORT}). \\
  \texttt{RESET\_CONTROL} & 0x010 & Software reset (\texttt{SOFT\_RST}) cho các khối và hàng đợi; chỉ có hiệu lực khi FSM ở \texttt{Idle}. \\
  \texttt{HC\_STATUS}     & 0x014 & Trạng thái của bộ điều khiển. \\
  \texttt{CMD\_QUEUE}     & 0x080 & Cổng ghi \CommandDescriptor{} vào \CommandFifo{}. \\
  \texttt{RESP}           & 0x084 & Cổng đọc \ResponseDescriptor{} từ \ResponseFifo{}. \\
  \texttt{PIO\_DATA\_PORT} & 0x088 & Cổng ghi dữ liệu TX vào \TxFifo{} và đọc dữ liệu RX từ \RxFifo{}. \\
  \texttt{QUEUE\_STATUS}  & 0x0B4 & Cờ đầy/rỗng của các hàng đợi. \\
  \texttt{T\_R}           & 0x32C & Thời gian cạnh lên SCL (I3C). \\
  \texttt{T\_F}           & 0x330 & Thời gian cạnh xuống SCL (I3C). \\
  \texttt{T\_SU\_DAT}     & 0x334 & Thời gian setup dữ liệu (I3C). \\
  \texttt{I2C\_T\_SU\_DAT} & 0x338 & Thời gian setup dữ liệu (\iic{}). \\
  \texttt{T\_HD\_DAT}     & 0x33C & Thời gian hold dữ liệu (I3C). \\
  \texttt{T\_HIGH}        & 0x340 & Thời gian SCL mức cao (I3C). \\
  \texttt{I2C\_T\_HIGH}   & 0x34C & Thời gian SCL mức cao (\iic{}). \\
  \texttt{T\_LOW}         & 0x350 & Thời gian SCL mức thấp (I3C). \\
  \texttt{T\_LOW\_OD}     & 0x354 & Thời gian SCL mức thấp ở \OpenDrain{}. \\
  \texttt{I2C\_T\_LOW}    & 0x358 & Thời gian SCL mức thấp (\iic{}). \\
  \texttt{T\_HD\_STA}     & 0x35C & Thời gian hold của START (I3C). \\
  \texttt{I2C\_T\_HD\_STA} & 0x360 & Thời gian hold của START (\iic{}). \\
  \texttt{T\_SU\_STA}     & 0x368 & Thời gian setup của \RepeatedSTART{} (I3C). \\
  \texttt{I2C\_T\_SU\_STA} & 0x36C & Thời gian setup của \RepeatedSTART{} (\iic{}). \\
  \texttt{T\_SU\_STO}     & 0x370 & Thời gian setup của STOP (I3C). \\
  \texttt{I2C\_T\_SU\_STO} & 0x374 & Thời gian setup của STOP (\iic{}). \\
  \texttt{T\_BUS\_FREE}   & 0x37C & Thời gian cho \BusFreeCondition{} (I3C). \\
  \texttt{I2C\_T\_BUF}    & 0x380 & Thời gian cho \BusFreeCondition{} (\iic{}). \\
  \texttt{DAT\_BASE}      & 0x400 & Địa chỉ nền của bảng DAT \DatEntryCount{} mục (mỗi mục \DatEntryWidth{}-bit). \\
  \bottomrule
\end{longtable}

Mỗi mục DAT (\texttt{dat\_entry\_t}) rộng \DatEntryWidth{}-bit, bố trí trường như sau:

\begin{longtable}{C{3.8cm} C{2.4cm} L{6.6cm}}
  \toprule
  \textbf{Trường} & \textbf{Bit} & \textbf{Ý nghĩa} \\
  \midrule
  \endhead
  \texttt{device}           & [31]    & Cờ phân loại: 1 = thiết bị \iic{} legacy, 0 = thiết bị I3C. \\
  (reserved)                & [30:23] & Luôn đọc về 0. \\
  \texttt{dynamic\_address} & [22:16] & \DynamicAddress{} I3C. \\
  (reserved)                & [15:7]  & Luôn đọc về 0. \\
  \texttt{static\_address}  & [6:0]   & \StaticAddress{} \iic{}. \\
  \bottomrule
\end{longtable}
